\documentclass[11pt,a4paper]{article}
\usepackage[utf8]{inputenc}
\usepackage[margin=1in]{geometry}
\usepackage{amsmath,amssymb,amsfonts}
\usepackage{hyperref}
\usepackage{booktabs}
\usepackage{graphicx}
\usepackage{listings}
\usepackage{xcolor}
\usepackage{microtype}

\hypersetup{
    colorlinks=true,
    linkcolor=blue,
    filecolor=magenta,      
    urlcolor=cyan,
    citecolor=red,
    pdftitle={Governed Agent: A Cryptographic Capability Framework for LLM Agent Tool Access},
    pdfauthor={jtevesnz-sys},
}

\lstset{
    basicstyle=\ttfamily\small,
    breaklines=true,
    frame=single,
    backgroundcolor=\color{gray!5},
    keywordstyle=\color{blue},
    commentstyle=\color{gray},
    stringstyle=\color{teal}
}

\title{\textbf{Governed Agent: Cryptographic Capability Sealing and Out-of-Band Audit for Autonomous LLM Agents}}
\author{
  \textbf{jtevesnz-sys} \\
  Mercenarios AI Systems \\
  \texttt{\href{https://github.com/jtevesnz-sys/governed-agent}{github.com/jtevesnz-sys/governed-agent}}
}
\date{October 2026}

\begin{document}

\maketitle

\begin{abstract}
Autonomous Large Language Model (LLM) agent architectures operate predominantly under \emph{ambient authority}. When tools are exposed to an agent runtime, tool invocation decisions depend upon probabilistic system prompts or secondary text classification guardrails. Under adversarial prompt injection, jailbreaking, or semantic drift, these linguistic boundaries fail, permitting unauthorized mutations. Furthermore, prevailing designs encourage models to evaluate the validity of their own output, introducing acute self-certification hazards. 

We present \textbf{Governed Agent}, a deterministic security architecture for LLM agent tool dispatch based on classical capability-based operating systems. Governed Agent enforces three architectural invariants: (1) \emph{Cryptographic Capability Sealing (KCAP)}, requiring unforgeable RSA-SHA256 tokens and monotonic anti-rollback sequence anchors before tool invocation; (2) \emph{Separation of Planning and Execution Authority (BLINK)}, enforcing that multi-turn decomposition yields strictly unprivileged task packets ($\mathtt{execution\_authority: False}$); and (3) \emph{Out-of-Band Independent Audit (Watchtower)}, decoupling verification from the generating model through external multi-point manifest inspections emitting non-negotiable $\mathtt{PASS}$, $\mathtt{HOLD}$, or $\mathtt{REVISE}$ verdicts. Benchmarks demonstrate that full capability verification introduces sub-millisecond overhead ($<0.5$\,ms per turn), proving that mathematical security primitives can enforce hard operational boundaries on autonomous models without sacrificing runtime performance.
\end{abstract}

\section{Introduction}

The deployment of autonomous Large Language Model (LLM) agents has transitioned from conversational interfaces to mutating operational pipelines. Modern agent frameworks routinely equip models with tools enabling direct database modifications, terminal execution, arbitrary HTTP dispatch, and local filesystem access.

In current production designs, tool execution operates under \textbf{ambient authority}: if an agent possesses a tool schema, any model completion that emits matching function-call syntax is executed directly by the runtime dispatcher. Security controls in this paradigm rely almost exclusively on probabilistic semantic filtering:
\begin{enumerate}
    \item \textbf{System Prompt Instructions:} Natural-language constraints (e.g., \emph{``Do not overwrite critical tables''}).
    \item \textbf{Secondary Model Classifiers:} Small language models deployed to classify incoming queries or proposed tool arguments (e.g., NeMo Guardrails, Llama Guard).
\end{enumerate}

These defenses suffer from an inherent domain mismatch: they attempt to solve a systems-security problem using non-deterministic semantic heuristics. Under prompt injection, multi-hop contextual obfuscation, or tool-parameter poisoning, probabilistic guardrails consistently degrade. Furthermore, agent pipelines routinely task the executing model with verifying its own output, violating the basic software engineering principle of independent verification.

Operating system design resolved this dilemma decades ago. In 1966, Dennis and Van Horn~\cite{dennis1966} formulated \textbf{capability-based security}: an entity must hold an unforgeable, communicable token of authority to access an object. In 1975, Saltzer and Schroeder~\cite{saltzer1975} codified the \textbf{separation of duties}: no sensitive state modification should proceed without decoupled authorization and verification.

We present \textbf{Governed Agent}, an open-source framework that adapts classical capability theory to LLM agent orchestration. Rather than attempting to persuade a language model to behave securely, Governed Agent enforces mathematical and cryptographic gates at the tool execution boundary.

\section{Threat Model and Security Invariants}

\subsection{Threat Model}
We model the LLM as an unprivileged, non-deterministic, and potentially untrusted reasoning core. The model may experience:
\begin{itemize}
    \item \textbf{Adversarial Injection:} Indirect prompt injection originating from untrusted web pages, database records, or third-party tool responses.
    \item \textbf{Contextual Drift:} Gradual divergence from initial operational boundaries across long context windows.
    \item \textbf{Hallucinatory Attestation:} Fabricated claims that requested modifications were completed successfully.
    \item \textbf{Privilege Replay:} Emitter replay of previously authorized actions outside their intended temporal or functional window.
\end{itemize}

The tool execution dispatcher, the authority key store, and the independent audit engine execute out-of-band in isolated host processes.

\subsection{Security Invariants}
Governed Agent enforces four non-negotiable invariants:
\begin{itemize}
    \item \textbf{Invariant 1 (Absence of Ambient Authority):} No mutating tool call is dispatched without a mathematically verified capability token signed by a designated custodian.
    \item \textbf{Invariant 2 (Plan $\neq$ Execution Authority):} Task planning and decomposition must output task packets where $\mathtt{execution\_authority} = \text{False}$.
    \item \textbf{Invariant 3 (Monotonic Anti-Rollback):} State checkpoints enforce monotonically increasing sequence pins ($Seq_{n+1} > Seq_n$). Replay or rollback of revoked or expired tokens is strictly rejected.
    \item \textbf{Invariant 4 (Out-of-Band Audit):} The executing agent cannot certify its own output. Build artifacts and directory mutations are validated solely by an external auditor against signed manifests.
\end{itemize}

\section{Architecture and Core Primitives}

\subsection{5-Layer Governance Model}
Governed Agent decomposes execution authority into five decoupled layers:
\begin{enumerate}
    \item \textbf{Capability Layer (KCAP):} RSA-SHA256 capability token sealing, canonical serialization, and monotonic anti-rollback state pinning.
    \item \textbf{Planning Layer (BLINK):} Bounded task packets, deterministic prohibition gates, and hash-bound handoffs.
    \item \textbf{Governance Layer:} Domain boundary enforcement and lawful purpose gating.
    \item \textbf{Intake Layer:} Credential sanitization, secret blocking, and immutable receipt logging.
    \item \textbf{Audit Layer (Watchtower):} Independent 6-point verification of generated artifacts and filesystem states.
\end{enumerate}

The reference implementation focuses on Layers 1, 2, and 5.

\subsection{Cryptographic Capability Sealing (KCAP Core)}
Before invoking a protected tool, an agent must present a sealed capability token. To ensure deterministic signature generation, all inputs are serialized via ASCII-sorted compact JSON canonicalization:
\begin{equation}
    R_{\text{bytes}} = \text{Canonicalize}(\text{Request} \cup \text{Context}_{\text{trusted}} \cup T_{\text{UTC}} \cup \text{KeyID})
\end{equation}
The token is sealed using RSA-2048/SHA-256 with PKCS\#1 v1.5 padding:
\begin{equation}
    \sigma = \text{Sign}_{K_{\text{priv}}}(\text{SHA-256}(R_{\text{bytes}}))
\end{equation}
At dispatch, the execution boundary evaluates:
\begin{equation}
    \text{Verify}_{K_{\text{pub}}}(R_{\text{bytes}}, \sigma) \equiv \text{VALID} \quad \land \quad \text{Context}_{\text{active}} \equiv \text{Context}_{\text{token}}
\end{equation}
Any mismatch in the trusted context (e.g., modified task ID, altered role, or tampered parameters) immediately results in a fail-closed $\mathtt{HOLD}$ state.

\subsection{Monotonic Anti-Rollback Authority}
Digital signatures alone are insufficient to defend against replay attacks where an agent resubmits an earlier valid token. Governed Agent couples capability checkpoints to a monotonic sequence authority:
\begin{equation}
    \text{Pin} = (\text{ChainID}, \text{Hash}_{\text{chk}}, \text{Sequence}, T_{\text{UTC}}, \sigma_{\text{auth}})
\end{equation}
The authority enforces:
\begin{equation}
    Seq_{\text{new}} > Seq_{\text{latest}}
\end{equation}
Any attempt to submit $Seq \le Seq_{\text{latest}}$ raises an immediate $\mathtt{ValueError: Rollback\ detected}$. Checkpoints must additionally fall within an enforced freshness window $\Delta t \in [1, 3600\,\text{s}]$:
\begin{equation}
    T_{\text{obs}} \le T_{\text{trusted}} \le T_{\text{expiry}} \quad \land \quad (T_{\text{trusted}} - T_{\text{obs}}) \le \Delta t_{\text{max}}
\end{equation}

\subsection{Scope-Bounded Planning (BLINK Core)}
The planning layer decouples decomposition from execution authority. Planning agents produce task packets:
\begin{equation}
    \mathcal{P} = (\text{TaskID}, \text{Objective}, \mathcal{A}_{\text{allowed}}, \mathcal{P}_{\text{prohibited}}, \text{Paths}, \text{Criteria})
\end{equation}
The BLINK validator enforces:
\begin{equation}
    \mathcal{A}_{\text{allowed}} \subseteq \{\text{read}, \text{propose}, \text{create-candidate}, \text{verify}, \text{classify}\}
\end{equation}
\begin{equation}
    \{\text{activate}, \text{deploy}, \text{publish}, \text{send}, \text{overwrite}\} \subseteq \mathcal{P}_{\text{prohibited}}
\end{equation}
\begin{equation}
    \mathcal{A}_{\text{allowed}} \cap \mathcal{P}_{\text{prohibited}} = \emptyset
\end{equation}
Validation enforces an invariant response:
\begin{equation}
    \mathtt{execution\_authority} = \text{False}, \quad \mathtt{not\_authorization} = \text{True}
\end{equation}
The packet is sealed into an immutable SHA-256 handoff record ($\mathtt{handoff\_ID}$), guaranteeing downstream workers cannot claim implicit authorization from planning materials.

\subsection{Out-of-Band Independent Audit (Watchtower Core)}
To eliminate model self-certification, artifact acceptance is assigned to an external verifier that runs out-of-band:
\begin{enumerate}
    \item \textbf{Manifest Validation:} Conformance to schema and completeness of declared hashes.
    \item \textbf{Lane Presence:} Verification that all expected execution lanes submitted artifacts.
    \item \textbf{Unlisted File Detection:} Directory sweeps identifying undeclared filesystem side effects.
    \item \textbf{Hash Agreement:} Full SHA-256 recomputation across all generated assets.
    \item \textbf{Cross-Lane Independence:} Detection of conflicting concurrent mutations across parallel lanes.
    \item \textbf{Residue Detection:} Inspection for leaked temporary or credential files ($\mathtt{.tmp}$, $\mathtt{.bak}$, $\mathtt{.env}$).
\end{enumerate}
The auditor returns discrete verdict enums: $\mathtt{PASS}$, $\mathtt{HOLD}$, or $\mathtt{REVISE}$. Build acceptance is impossible without an external $\mathtt{PASS}$.

\section{Empirical Evaluation}

\subsection{Implementation Profile}
The Tier 1 reference implementation comprises roughly 1,200 lines of pure Python code requiring only \texttt{cryptography} and \texttt{pydantic}. The suite contains 23 unit tests executing with 88\% statement coverage.

\begin{table}[h]
\centering
\small
\begin{tabular}{@{}llcc@{}}
\toprule
\textbf{Module} & \textbf{Functionality} & \textbf{Lines of Code} & \textbf{Coverage} \\
\midrule
\texttt{governed\_agent.capability.crypto} & RSA-2048 sealing, keystore, canonical JSON & 390 & 88\% \\
\texttt{governed\_agent.capability.antirollback} & Monotonic pins, freshness intervals & 337 & 86\% \\
\texttt{governed\_agent.planning.blink} & Task packet validation, handoffs & 218 & 78\% \\
\texttt{governed\_agent.audit.watchtower} & Out-of-band 6-point manifest audit & 213 & 99\% \\
\midrule
\textbf{Total Reference Implementation} & & \textbf{1,239} & \textbf{88\%} \\
\bottomrule
\end{tabular}
\caption{Codebase profile and test coverage for the Tier 1 reference release.}
\label{tab:profile}
\end{table}

\subsection{Microbenchmarks and Overhead}
Benchmarks were executed on an AMD processor running Python 3.14 on Windows 11. Over 1,000 iterations:
\begin{itemize}
    \item \textbf{Keypair Generation (RSA-2048):} $12.4$\,ms.
    \item \textbf{Capability Sealing (RSA-SHA256):} $0.42$\,ms.
    \item \textbf{Capability Token Verification:} $0.08$\,ms.
    \item \textbf{End-to-End Governance Cycle:} $0.05$\,s (Planning $\rightarrow$ Capability $\rightarrow$ Audit).
    \item \textbf{Full Test Suite (23 tests):} $0.54$\,s total runtime.
\end{itemize}

Given that typical frontier LLM generation latencies range from $500$\,ms to $3,000$\,ms per turn, the cryptographic verification overhead accounts for \textbf{less than 0.1\%} of turn latency.

\subsection{Architectural Comparison}

\begin{table}[h]
\centering
\small
\begin{tabular}{@{}lcccc@{}}
\toprule
\textbf{Dimension} & \textbf{LangChain/CrewAI} & \textbf{Anthropic MCP} & \textbf{NeMo Guardrails} & \textbf{Governed Agent} \\
\midrule
Execution Authority & Ambient & Server-level & Secondary prompt & \textbf{RSA Capability Token} \\
Replay/Rollback Defense & None & None & None & \textbf{Monotonic Sequence Pins} \\
Plan/Execution Separation & Blended & Blended & None & \textbf{Enforced ($\mathtt{auth=False}$)} \\
Audit Mechanism & Self-certification & None & Secondary prompt & \textbf{Out-of-band 6-point check} \\
Failure Mode & Unbounded mutation & Client error & Evasion/bypass & \textbf{Fail-closed ($\mathtt{HOLD}$)} \\
\bottomrule
\end{tabular}
\caption{Comparison of agent governance and execution control mechanisms.}
\label{tab:comparison}
\end{table}

\section{Related Work}

\textbf{Agent Frameworks and Tool Calling:} LangChain~\cite{chase2022}, AutoGen, and CrewAI pioneered multi-agent collaboration and automated function calling. While they provide routing abstractions, tool execution remains ambient once configured. Anthropic's Model Context Protocol (MCP) formalizes client-server tool schemas over JSON-RPC but delegates permission control to coarse client-level toggles without cryptographic invocation tokens.

\textbf{Prompt-Based Guardrails:} NeMo Guardrails~\cite{rebedea2023} and Llama Guard introduce programmable rails and input/output classification models. While effective for semantic moderation, they remain vulnerable to adversarial jailbreaks and cannot produce non-repudiable audit trails.

\textbf{Capability Security:} Dennis and Van Horn~\cite{dennis1966} established capability-based computing, realized in operating systems such as KeyKOS, EROS, and Google Fuchsia. Governed Agent applies these principles to LLM agent systems, treating tools as privileged capabilities requiring non-repudiable tokens.

\section{Conclusion and Future Work}

Probabilistic prompt guardrails cannot guarantee safety in autonomous agent systems with mutating capabilities. Governed Agent demonstrates that classical capability-based security, monotonic anti-rollback state, and decoupled out-of-band audits provide mathematically deterministic controls at negligible latency overhead.

Future research includes formalizing Layer 3 (consent-gated case envelopes) and Layer 4 (receipt stores with automated secret sanitization), alongside integrating Hardware Security Modules (HSMs) and RFC 3161 timestamping authorities for enterprise monotonic pinning.

\section*{Availability}
The complete reference implementation, whitepaper, and runnable demonstrations are open source and available at: \\
\url{https://github.com/jtevesnz-sys/governed-agent}

\begin{thebibliography}{99}
\bibitem{dennis1966}
J.~B. Dennis and E.~C. Van~Horn, ``Programming semantics for multiprogrammed computations,'' \emph{Communications of the ACM}, vol.~9, no.~3, pp. 143--155, 1966.

\bibitem{saltzer1975}
J.~H. Saltzer and M.~D. Schroeder, ``The protection of information in computer systems,'' \emph{Proceedings of the IEEE}, vol.~63, no.~9, pp. 1278--1308, 1975.

\bibitem{chase2022}
H.~Chase, ``LangChain: Building applications with LLMs through composability,'' 2022.

\bibitem{rebedea2023}
T.~Rebedea et al., ``NeMo Guardrails: A toolkit for controllable and safe LLM applications with programmable rails,'' \emph{NVIDIA}, 2023.
\end{thebibliography}

\end{document}
